\scanpage{66}
% ll. 1–8
\deleted{Donc le choix} (ce qui ne \unsure{;} peut être le cas que pour un seul des entiers $i=0,1$) alors la donnée de $(\lambda_0,\lambda_1,K)$ ne détermine pas la classe d'iso \unsure{du|de l'|des} hyperpolygône pythagoricien, i.e. ne détermine pas $\zeta$

% ll. 9–12
\uline{En résumé}

a) Pour \diagram{\begin{tikzcd}[ampersand replacement=\&,row sep=.6em,column sep=.18em]
(\lambda_0,\arrow[d,leftright] \& \lambda_1,\arrow[d] \& K)\\
K_0 \& K_1 \&
\end{tikzcd}} admissible \inserted{(du pt de vue projectif)}, \uuline{si} $K_0,K_1\ne-2$ i.e. $\lambda_0,\lambda_1\ne-1$ i.e. $\nu_0,\nu_1\ne2$ il y a un \inserted{et un} seul hyperpolygône associé (mod iso) i.e. un et un seul $\zeta$ correspondant. Si $K=-2$, \inserted{i.e. $\nu=2$,} il est donc donné \deleted{par} \inserted{comme l'unique solution de} l'équation en $\zeta$ ; si $K\ne\unsure{-}2$ parmi les 2 solutions de cette équation, une seule convient pour \inserted{définir} un hyperpolygône pythagoricien. [Comment les distinguer ? Si on travaille dans $\mathbb C$, on peut songer à la condition $0<\zeta<1$ --- ]

% ll. 19–25
b) Par contre si $K_0$ ou $K_1=-2$ i.e. $\lambda_0$ ou $\lambda_1=-1$, \inserted{i.e. $\nu_0$ ou $\nu_{\unsure{1|i}}=2$} alors alors pour $(\lambda_0,\lambda_1,K)$ admissible \inserted{(projectivem.)} donné, il y a \uline{deux} hyperpolygônes non iso qui lui correspondent, correspondant donc aux 2 solutions distinctes $\zeta$ de (E) [NB on sait \inserted{par ailleurs} que l'on aura alors $K=-2$ i.e. que l'on aura bien 2 racines distinctes].

% ll. 26–27
Il n'est donc pas vrai que $\nu_0,\nu_1,\nu$ déterminent la classe d'iso. d'un hyperpolygône, \unsure{$=$|même|---} la.
% archive sheet number: 65
